\section{An angular projection gain}\label{ang:section}

This section proves the finite-scale incidence estimate used to improve
the angular exponent in odd dimensions. Its geometric content is a bound
on incidences whose fiber at each line direction is covered by a small
number of parallel tubes. The planar discretized Furstenberg theorem is
the only deep input.

All sets in this section lie in fixed bounded regions. Directions are unit
vectors, with antipodal vectors identified when discussing lines. A finite
number of projective charts removes this distinction. For a bounded set
$A$, write $|A|_r$ for its covering number by balls of radius $r$. A tube has
bounded length and the stated perpendicular width; enlarging its length by
a fixed constant is always permitted. We use the ball-bound and subpower
conventions of Section~\ref{sec:conventions}. Bounds for radii in a fixed
larger bounded interval follow by changing constants. In particular,
incidence mass at least subpower means a lower bound $\delta^{o(1)}$.

\subsection{The gain and its two slicing regimes}

\begin{theorem}[Linear projection gain]\label{ang:linear-gain}
Fix $k\ge2$ and
\[
 k-1\le q\le k-\tfrac12<t<k.
\]
Fix a tube-width constant $C$. In the case $q=k-1$, also fix
$\zeta>0$. There exist $\epsilon,c>0$ such that, for all sufficiently
small $\delta$, the following holds. Let $\mu$ be a bounded point
subprobability in $\R^{k+1}$ and $\nu$ a direction subprobability on $S^k$
with respective $t$- and $q$-ball bounds of loss $\delta^{-\epsilon}$
down to $\delta$. At $q=k-1$ assume in addition
\begin{equation}\label{ang:endpoint-plate}
 \nu\{w:\dist(w,H)\le r\}\le\delta^{-\epsilon}r^\zeta
       \qquad(\delta\le r\le1)
\end{equation}
for every linear hyperplane $H$. Let $E\subset\R^{k+1}\times S^k$ be a
measurable incidence relation, and write $E_w=\{x:(x,w)\in E\}$.
If each direction fiber $E_w$ is covered by at most
$\delta^{-q-\epsilon}$ tubes parallel to $w$ of width $C\delta$, then
\[
 (\mu\otimes\nu)(E)\le\delta^c.
\]
All constants are uniform over the measures and the incidence relations
with these fixed parameters.
\end{theorem}

A cover by $\delta^{-q-\epsilon}$ tubes means that the perpendicular
projection of a direction fiber is covered by that many balls of radius
comparable to $\delta$. The theorem says that such covers cannot capture
more than power-small incidence between the two nonconcentrated laws.

The two cases in its proof are distinguished by what remains after
slicing. If $q>k-1$, slicing by codimension $h=k-1$ leaves a plane and
exponents
\[
 b=q-h\in(0,\tfrac12],\qquad f=t-h\in(\tfrac12,1),\qquad b<f.
\]
The planar incidence estimate applies in this range. At $q=k-1$, the
same slicing would leave direction exponent zero. Instead we use
$h=k-2$, leaving three-dimensional slices with direction exponent one
and point exponent greater than $3/2$. The endpoint argument finds three
transverse directions in one such slice. A weak radial estimate with
exponent $\ell>k-2$ controls the triples that are almost dependent.

\subsection{The planar incidence input and its measure formulation}

The metric on bounded affine line parameters may be taken to be the sum
of the distance between the associated orthogonal projection operators
and the distance between the closest-to-origin points on the lines.
Other smooth bounded-chart metrics are equivalent. A nonempty bounded
set $P$ of points or line parameters is a $(\delta,s,K)$-set if
\[
 |P\cap B(a,r)|_\delta\le K r^s |P|_\delta
       \qquad(\delta\le r\le1).
\]
The normalization by $|P|_\delta$ is part of the defining bound.

\begin{theorem}[Planar discretized Furstenberg estimate]
\label{ang:furstenberg}
Let $0<b<1$ and $b<f<2$. There is $\epsilon_0>0$ such that the
following holds for sufficiently small $\delta$. Suppose that $P$ is
a nonempty bounded $(\delta,f,\delta^{-\epsilon_0})$-set and that
each $x\in P$ has a nonempty family $\mathcal T_x$ of incident
$\delta$-tubes whose line parameters form a
$(\delta,b,\delta^{-\epsilon_0})$-set. Then
\[
 \left|\bigcup_{x\in P}\mathcal T_x\right|_\delta
       \ge\delta^{-2b-\epsilon_0}.
\]
The exponent may be decreased to absorb fixed bounded charts and fixed
width multiples.
\end{theorem}

This is the tube formulation in
\cite[Definitions~2.5--2.6 and Theorem~2.7]{OSW2024}, following the
discretized theorem of \cite[Theorem~1.3]{OS2023}. In the application
below, $f<1$; that range is included in the theorem.

The measure formulation below keeps the labels being sampled separate
from their geometric coordinates. Several labels may have the same
coordinate and still carry different incidence fibers.

\begin{lemma}[Planar measure incidence obstruction]
\label{ang:planar-measures}
Fix $0<b<1$ and $b<f<2$. There is no sequence $\delta\to0$ with the
following data. Let $\mu,\nu$ be subprobabilities on Borel label sets
$X,W$ in Euclidean spaces, and let
\[
 p:X\longrightarrow\R^2,\qquad v:W\longrightarrow S^1
\]
be measurable coordinate maps with fixed bounded images. Suppose that
$p_\#\mu$ and $v_\#\nu$ have respective $f$- and $b$-ball bounds
of loss $\delta^{-o(1)}$ down to $\delta$. Let $E\subset X\times W$ be a measurable
incidence relation with
\[
 (\mu\otimes\nu)(E)\ge\delta^{o(1)}.
\]
For each original direction label $\omega\in W$, put
$E_\omega=\{\xi\in X:(\xi,\omega)\in E\}$. Assume that the set of
point coordinates $p(E_\omega)$ is covered by at most
$\delta^{-b-o(1)}$ strips of width $C\delta$ parallel to $v(\omega)$.
\end{lemma}

\begin{proof}
Restrict the coordinate images to fixed bounded charts and normalize the
two laws. Their masses are at least the incidence mass, so this costs only
subpower factors. Let $I$ be the resulting incidence mass, sample
$n=\lceil\delta^{-b}\rceil$ independent direction labels
$\omega_1,\ldots,\omega_n$ from $\nu$, and write $v_i=v(\omega_i)$.
There is a subpower quantity $T\to\infty$ such that, with failure
probability smaller than any prescribed power of $\delta$, the sample
satisfies
\begin{equation}\label{ang:sample-caps}
 \#\{i:v_i\in B(v_0,r)\}\le Tnr^b
       \qquad(\delta\le r\le1).
\end{equation}
To verify this, use a fixed collection of dyadic grids in the direction
chart. There are only polynomially many relevant cells. Their expected
counts are at most $Knr^b$, where $K=\delta^{-o(1)}$, and
$nr^b\ge1$. Choose $T$ larger than a fixed large multiple of
$K+\log(1/\delta)$. For a binomial count $Y$, the bound
\[
 \Pr(Y\ge j)\le (e\mathbb E Y/j)^j
\]
and a union bound prove Equation~\eqref{ang:sample-caps}, including
arbitrary balls by fixed grid comparison. Increasing $T$ by another
subpower factor makes the failure probability smaller than $I/4$.

The expected empirical incidence mass equals $I$. Since it is bounded
by one, there exists a sample satisfying Equation~\eqref{ang:sample-caps}
and having empirical incidence mass at least $I/2$. The set $X_0$ of
point labels incident to at least $In/4$ sampled direction labels has
$\mu$-mass at least $I/4$: its complement contributes at most $I/4$
to the empirical incidence integral. Normalize $\mu|_{X_0}$ and sample
$m=\lceil\delta^{-f}\rceil$ independent point labels. If their coordinates
are $x_1,\ldots,x_m$, the same argument gives a sample with
\[
 \#\{j:x_j\in B(x,r)\}\le T'mr^f,
 \qquad T'=\delta^{-o(1)},\qquad \delta\le r\le1.
\]
In particular, the multiplicities in $\delta$-balls are subpower.
Color a spatial $\delta$-grid with finitely many colors so that cells
with the same color are $\delta$-separated. Choose a color with the most
occupied cells and one sampled label from each of them. Their
point coordinates form a separated set $P$ with
$|P|_\delta\ge m\delta^{o(1)}$. The displayed local count makes $P$
a $(\delta,f,\delta^{-o(1)})$-set. Each $x\in P$ retains the list of at
least $In/4$ good sampled direction labels belonging to its chosen point
label.

For $x\in P$, take the lines exactly through $x$ in the coordinates
$v_i$ of those listed labels, and let $\mathcal T_x$ be their
$\delta$-tubes. Equation~\eqref{ang:sample-caps} bounds the multiplicity
in a $\delta$-ball of directions by $CT$, and the count in an $r$-ball
by $Tnr^b$. Consequently
\[
 |\mathcal T_x|_\delta\ge cIn/T,
 \qquad
 |\mathcal T_x\cap B(\ell,r)|_\delta\le CTnr^b.
\]
Here the direction component of the line metric supplies the lower
distance comparison; the reverse comparison holds because $x$ is bounded.
The normalized loss is therefore at most $CT^2/I=\delta^{-o(1)}$.
In particular, $\mathcal T_x$ is a nonempty
$(\delta,b,\delta^{-o(1)})$-set of line parameters even when different
sampled labels have the same direction coordinate.

For each sampled label $\omega_i$, its prescribed strips contain the
coordinates of all its incident sampled point labels. Moving the central
line of such a strip to pass through one of these points changes its
intercept by $O(\delta)$. Thus each original strip accounts for only a
bounded number of $\delta$-cells in line-parameter space. Altogether the
union of our line families has covering count at most
\[
 n\delta^{-b-o(1)}=\delta^{-2b-o(1)}.
\]
This contradicts Theorem~\ref{ang:furstenberg} for small $\delta$.
The count treats each original sampled direction label separately; it
never asks for one common fiber cover after coincident coordinates are
identified.
\end{proof}

\subsection{Slicing with controlled losses}

We state the needed estimates for Haar-random subspaces before applying
them. Haar measure is normalized to have total mass one.

\begin{lemma}[Haar band estimates]\label{ang:haar}
Let $0\le h\le n-1$, and let $V$ be a Haar-random codimension-$h$
linear subspace of $\R^n$.
For a unit vector $w$,
\[
 \Pr\{\dist(w,V)\le\delta\}\asymp\delta^h.
\]
If $j\le n-h$ unit vectors $w_1,\ldots,w_j$ have singular values
$d_1\ge\cdots\ge d_j\ge0$, then
\begin{equation}\label{ang:haar-multiple}
 \delta^{-jh}\Pr\{\dist(w_i,V)\le\delta\text{ for all }i\}
       \le C_{n,j}\prod_{i=1}^j(d_i+\delta)^{-h}.
\end{equation}
When $h=0$, the probability and the factors with exponent $h$ are one.
\end{lemma}

\begin{proof}
The one-vector estimate follows from spherical coordinates. For the
second estimate, singular-value decomposition implies that the band
conditions constrain the projections of orthonormal singular directions
$u_i$ by $C_j\delta/d_i$. Keep only the constraints whose right-hand
sides are less than a sufficiently small dimensional constant, and
order these right-hand sides increasingly. By rotational invariance,
one can fix $V^\perp$ and reveal successive vectors of a random
orthonormal frame. Conditional on the earlier constrained vectors,
projection onto $V^\perp$, restricted to their orthogonal complement,
has rank $h$ and nonzero singular values bounded below. Indeed the sum
of squared lengths of the earlier projections is less than $1/2$ if
the dimensional constant was chosen small enough. The next vector is
uniform on the unit sphere of that complement. Its probability of
having projection of length at most $r$ is at most $C_n r^h$.
The ambient dimension at this last conditional step is at least
$n-j+1\ge h+1$, including the limiting case of a one-dimensional
kernel. Multiplying the conditional estimates proves
\[
 \Pr\{w_i\text{ all in the band}\}
       \le C_{n,j}\prod_i\min(1,\delta/d_i)^h,
\]
which is equivalent to Equation~\eqref{ang:haar-multiple}.
\end{proof}

Shmerkin and Wang \cite[Sections~6.1--6.3]{SW2025} use normalized bands
to define slices and combine sliced with projected thin-tube bounds.
Here we keep bands at width $\delta$ and control the retained incidence
mass, ball bounds, and active tubes for each original direction label.

\begin{lemma}[Selection of slices]\label{ang:slice-selection}
Let $0\le h\le n-2$. Let point and direction subprobabilities $\mu,\nu$ in $\R^n$ have
$t$- and $q$-ball bounds with subpower losses, where $t,q>h$.
Assume the directions lie in a fixed small projective chart. Let an
incidence relation $E$ have mass at least subpower, and suppose each
direction fiber is covered by at most $\delta^{-q-o(1)}$ tubes of
width $C\delta$.

For Haar-random codimension-$h$ subspaces $V$ and
$z\in V^\perp$, set
\begin{align*}
 d\nu_V(w)&=\delta^{-h}\one_{\{|P_{V^\perp}w|\le\delta\}}\,d\nu(w),\\
 d\mu_{V,z}(x)&=\delta^{-h}
       \one_{\{|P_{V^\perp}x-z|\le\delta\}}\,d\mu(x).
\end{align*}
Integrate $z$ over a fixed sufficiently large ball. There are families of restrictions
$\widetilde\nu_V\le\nu_V$, $\widetilde\mu_{V,z}\le\mu_{V,z}$,
and relations $\widetilde E_{V,z}\subset E$, measurable in the displayed
parameters and restricted to the two bands defining $\mu_{V,z}$ and
$\nu_V$, such that
\[
 \mathbb E_V\int
 (\widetilde\mu_{V,z}\otimes\widetilde\nu_V)(\widetilde E_{V,z})\,dz
 \ge\delta^{o(1)},
\]
and the following hold:
\begin{enumerate}
\item The restricted measures have masses and ball losses at most
$\delta^{-o(1)}$, with exponents $q-h$ and $t-h$, respectively.
\item Each remaining direction fiber is covered by at most
$\delta^{-(q-h)-o(1)}$ width-$C'\delta$ tubes.
\item Put
\[
 p_{V,z}(x)=P_Vx,\qquad v_V(w)=P_Vw/|P_Vw|
 \quad\text{for }|P_{V^\perp}w|\le\delta.
\]
The first map is projection to $V+z$ followed by translation to $V$;
the second is well-defined on the direction band for small $\delta$.
The coordinate pushforwards of the restricted measures under these maps
satisfy the ball bounds in item~1, with fixed constant changes. For every
original direction label $w$ in that band, the coordinates
$p_{V,z}(\{x:(x,w)\in\widetilde E_{V,z}\})$ have the tube cover in
item~2 with direction $v_V(w)$. This cover remains attached to $w$ even
when two projected direction coordinates coincide.
\end{enumerate}
The direction restriction $\widetilde\nu_V$ depends on $V$ only, and its
mass has an a priori subpower upper bound valid for every $V$.
\end{lemma}

\begin{proof}
We first choose the fiber covers measurably, enlarging their width by a
fixed dimensional factor while preserving each actual direction $w$.
Let $M_\delta=\delta^{-q-o(1)}$ be their common size bound. Use a finite
spatial $\delta$-grid in a fixed bounded region containing one point on
every axis whose tube meets the source region. Rounding such a point
to the grid and keeping the direction $w$ enlarges the width by only
$C_n\delta$. Candidate covers are the finitely many lists of at most
$\lfloor M_\delta\rfloor$ grid anchors. For each candidate $J$, the
function
\[
 w\longmapsto
 \mu\left(E_w\setminus\bigcup_{a\in J}T(w,a)\right)
\]
is measurable by Fubini, where $T(w,a)$ is the enlarged tube parallel
to $w$ through $a$. Choose the first candidate for which this function
is zero. Coverage up to a $\mu$-null set suffices, since every point
slice $\mu_{V,z}$ is absolutely continuous with respect to $\mu$.
For a relation measurable only in the completed product sigma-algebra,
first choose a Borel subset agreeing with the relation modulo a
product-null set, and discard the $\nu$-null set of
directions on which its fiber differs by positive $\mu$-mass. Every
direction slice is absolutely continuous with respect to $\nu$, so
this replacement also changes no sliced incidence.
Writing $J(w)$ for the selected candidate, now replace $E$ by
\[
 E\cap\left\{(x,w):x\in\bigcup_{a\in J(w)}T(w,a)\right\},
\]
with the exceptional null directions removed. This loses zero original
or sliced incidence and makes the selected covers contain every fiber
of the restricted relation.

Intersecting the relation with the two bands does not change its sliced
incidence. The expected sliced incidence is comparable from below to the
original one. In fact, integrating the point-band indicator in $z$ gives a
constant times $\delta^h$, and the one-vector Haar estimate supplies
another factor comparable to $\delta^h$. These cancel the two slice
normalizations.

Write $M(V)=\|\nu_V\|$ and $L(V,z)=\|\mu_{V,z}\|$.
Uniformly in $V$,
\begin{equation}\label{ang:point-slice-total}
 \int L(V,z)\,dz\le C.
\end{equation}
The two-vector estimate in a projective chart gives the direction-pair
kernel
\begin{equation}\label{ang:direction-pair-kernel}
 \delta^{-2h}\Pr\{w,w'\text{ both within }\delta\text{ of }V\}
       \le C(|w-w'|+\delta)^{-h}.
\end{equation}
For points, integrate both band indicators in $z$ first. This yields
at most $C\delta^h$ times the indicator
$|P_{V^\perp}(x-x')|\le2\delta$. Haar averaging therefore gives
the analogous normalized kernel $C(|x-x'|+\delta)^{-h}$.
Dyadic annulus summation, using $q,t>h$, proves
\begin{align}
 \mathbb E_V M(V)^2+\mathbb E_V\int L(V,z)^2\,dz
       &\le\delta^{-o(1)},\label{ang:slice-second}\\
 \mathbb E_V\iint_{|w-w'|\le r}d\nu_V(w)d\nu_V(w')
       &\le\delta^{-o(1)}r^{q-h},\label{ang:direction-close}\\
 \mathbb E_V\int\iint_{|x-x'|\le r}
             d\mu_{V,z}(x)d\mu_{V,z}(x')\,dz
       &\le\delta^{-o(1)}r^{t-h}\label{ang:point-close}
\end{align}
for $\delta\le r\le1$. For example, the shells below $r$ form a
geometric sum with successive exponents $q-h>0$ in
Equation~\eqref{ang:direction-close}; the finest shell uses the
original ball estimate at radius $\delta$.

The truncations below first bound direction-slice mass, then enforce
both ball bounds, and finally bound the active tube count. Take a subpower number
$K\to\infty$ absorbing all constants, all input losses, the dyadic
radius sums in Equations~\eqref{ang:direction-close} and
\eqref{ang:point-close}, and the inverse original incidence mass.
Enlarge it so the sliced incidence is at least $K^{-1}$ and the
right-hand side of Equation~\eqref{ang:slice-second} is at most $K$.
First remove planes with $M(V)>K^3$. By
Equations~\eqref{ang:point-slice-total} and \eqref{ang:slice-second},
their incidence contribution is at most $CK^{-2}$.

In the remaining point slices, remove points at which
\[
 \mu_{V,z}(B(x,Cr))>K^6r^{t-h}
\]
at some dyadic radius $r\in[\delta,1]$. Markov's inequality and
Equation~\eqref{ang:point-close} bound the expected deleted point mass
by $CK^{-5}$. Since $M(V)\le K^3$, its incidence cost is at most
$CK^{-2}$. The restricted point measure has the asserted ball bound
at every center: any ball meeting the retained set can be enlarged
by a fixed factor around a retained point. It also has mass $O(K^6)$,
by covering the bounded support with a fixed number of unit balls.

Similarly remove direction-slice points for which
$\nu_V(B(w,Cr))>K^3r^{q-h}$ at some dyadic radius. Their expected
mass is at most $CK^{-2}$ by
Equation~\eqref{ang:direction-close}; Equation~\eqref{ang:point-slice-total}
bounds the corresponding integrated incidence loss by the same quantity.
No independence between direction mass and point mass is used here.

Finally, for a direction $w$ in the band about $V$, let $J(w)$ be its
selected list of grid anchors. A bounded tube through $a$, parallel to
$w$, can meet the source region in the band about $V+z$ only if
$|P_{V^\perp}a-z|\le C'\delta$: its transverse projection within the
bounded source region has diameter $O(\delta)$. Thus the measurable
upper count
\[
 N(V,z,w)=\sum_{a\in J(w)}
       \one_{\{|P_{V^\perp}a-z|\le C'\delta\}}
\]
dominates the number of active covering tubes and satisfies
\[
 \int N(V,z,w)\,dz
       \le C M_\delta\delta^h\le K\delta^{h-q}.
\]
Delete direction/$z$ pairs for which
$N(V,z,w)>K^{12}\delta^{h-q}$. Their $z$-volume for each $V,w$
is at most $K^{-11}$. The truncated point mass is $O(K^6)$ and
the truncated direction mass is at most $K^3$, so this costs at most
$CK^{-2}$ in integrated incidence. Delete the corresponding incidences
but keep the marginal restrictions already constructed. At large $K$
the surviving integrated incidence is at least $(2K)^{-1}$.

Projection changes each direction by $O(\delta)$. The points of one
band whose coordinates $p_{V,z}(x)$ lie in a ball of radius $r\ge\delta$
lie in a ball of radius $Cr$ in the original space; the same assertion
holds for directions after unit normalization. These preimage bounds give
the asserted ball bounds on the coordinate pushforwards. For each original
label $w$, every active tube in its own selected list becomes a tube of
width $C'\delta$ parallel to $v_V(w)$. This proves item~3 for each
original direction label.
\end{proof}

\subsection{Endpoint preparation: transverse tubes and weak radial estimates}

We first record the elementary counting fact used twice below.

\begin{lemma}[Transverse tube counting]\label{ang:loomis-whitney}
Let $v_1,\ldots,v_m$ be unit vectors in $\R^m$, $m\ge2$, and suppose
$|\det(v_1,\ldots,v_m)|\ge\Delta>0$. If a bounded set $A$ is covered,
for each $i$, by $M_i$ tubes of width $C\delta$ parallel to $v_i$, then
\[
 |A|_{C_m\delta}\le C_{m,C}\Delta^{-m}
       \prod_{i=1}^m M_i^{1/(m-1)}.
\]
Changing the fixed constant in the covering radius changes only the
constant in this inequality.
\end{lemma}

\begin{proof}
Let $L$ be the matrix with columns $v_i$. Then $\|L\|\le C_m$ and
$\|L^{-1}\|\le C_m\Delta^{-1}$. Let $S\subset\Z^m$ index the
$\delta$-grid cubes meeting $L^{-1}A$. The image of such a cube has
diameter at most $C_m\delta$. Each original tube parallel to $v_i$
becomes a tube parallel to the $i$th coordinate axis, of width at most
$C_{m,C}\delta\Delta^{-1}$. Therefore
\[
 |\operatorname{pr}_{\widehat i}S|
       \le C_{m,C}\Delta^{-(m-1)}M_i.
\]
The discrete Loomis--Whitney inequality \cite{LoomisWhitney1949}
$|S|^{m-1}\le\prod_i|\operatorname{pr}_{\widehat i}S|$ proves the
claim. For completeness, if $X$ is uniform on $S$, the entropy chain rule
gives
\[
 H(X_{\widehat i})\ge
 \sum_{j\ne i}H(X_j\mid X_1,\ldots,X_{j-1}).
\]
Summing in $i$ yields $(m-1)\log|S|\le
\sum_i\log|\operatorname{pr}_{\widehat i}S|$, as required.
\end{proof}

For $x\ne z$, put $\pi_x(z)=(z-x)/|z-x|$. A restriction of the
product measure to a pair set is not normalized in the next statement.

\begin{lemma}[Weak radial estimate]\label{ang:weak-radial}
Fix $k\ge2$, $\zeta>0$, and
\[
 k-2<\ell<\frac{(k-1)^2}{k}.
\]
For every $e>0$ there exist $\epsilon,c>0$ with the following property
for sufficiently small $\delta$. Suppose that a probability $\lambda$
in a bounded subset of $\R^k$ satisfies
\begin{align*}
 \lambda(B(x,r))&\le\delta^{-\epsilon}r^{k-1},\\
 \lambda(H^{(r)})&\le\delta^{-\epsilon}r^\zeta
 \qquad(\delta\le r\le1)
\end{align*}
for every affine hyperplane $H$, where $H^{(r)}$ is its $r$-neighborhood.
There is a measurable pair set $G$ such that
\[
 (\lambda\otimes\lambda)(G^c)\lesssim\delta^c,
 \qquad
 \lambda\{z:(x,z)\in G,\ \pi_x(z)\in B(\theta,r)\}
       \le\delta^{-e}r^\ell
\]
for every $x$, every spherical cap center $\theta$, and
$\delta\le r\le1$. The diagonal is omitted from $G$. The conclusion
also holds for balls in projective direction space, with a fixed change
of constant, by taking the two antipodal caps.
\end{lemma}

\begin{proof}
We first prove a single-scale assertion at $u=2^{-N}$, with reciprocal
scale $R=2^N$: for any fixed $\ell_1<(k-1)^2/k$, sufficiently small
fixed losses in the hypotheses permit deletion of $O(R^{-c_1})$ pair
mass so that every remaining fiber has cap mass at most $R^{-\ell_1}$
at radius $u$. All choices made below are fixed independently of $R$.

Apply Lemma~\ref{reg:regularization} to the source measure with a small
parameter $\tau>0$. It suffices to treat one normalized source class
$\rho$, retaining the original pin law $\lambda$. The discarded source
mass is a fixed negative power of $R$, and the classes will be recombined
with their original weights. Write $m$ for the class profile. Its
comparison errors are $R^{C_k\tau}$, it is $k$-Lipschitz, and
\[
 m(1)\ge k-1-C_k\tau-\epsilon.
\]
The constant absorbs the normalization of a class of mass at least
$R^{-\tau}$. At depths $a<b$, an occupied depth-$b$ child of an
occupied depth-$a$ cell $Q$ has conditional $\rho$-mass at most
\begin{equation}\label{ang:conditional-child}
 R^{C_k\tau}D^{-1},\qquad
 D=R^{m(b/N)-m(a/N)}.
\end{equation}

Divide the depth interval into $J$ nearly equal blocks, with $J$ a
large fixed integer, and ignore the first two blocks. In every later
block $[a,b]$, remove the pairs whose pin is within $R^{-1/(4J)}$ of
the center $z_Q$ of the source cell $Q$ at depth $a$. Summing with the
source-cell weights shows that this costs at most
$C_JR^{-(k-1)/(4J)+\epsilon}$ pair mass. For the other pairs the center
direction
\[
 v(x,Q)=\frac{z_Q-x}{|z_Q-x|}
\]
differs from the direction to any point of $Q$ by
$O(2^{-a}R^{1/(4J)})$. In particular the transverse displacement inside
$Q$ caused by freezing this direction is
\begin{equation}\label{ang:freezing-error}
 O(2^{-2a}R^{1/(4J)})=o(2^{-b});
\end{equation}
here $a\ge2N/J+O(1)$ and $b-a=N/J+O(1)$.

Fix a small $e'>0$. A pair $(x,z)$, $z\in Q$, is troublesome if the
full conditional source mass in the width-$C2^{-b}$ tube about the
line through $x,z$ exceeds
\begin{equation}\label{ang:local-fraction}
 F=R^{e'}D^{-(k-1)/k}.
\end{equation}
The fixed width $C$ may be chosen as large as any subsequent application
requires. If $F\ge1$, there are no troublesome pairs. Otherwise,
Equation~\eqref{ang:freezing-error} and a grid in $v(x,Q)^\perp$ show
that all troublesome sources for this pin lie in at most $C/F$
parallel tubes of width $C'2^{-b}$. Indeed, each relevant projected grid
cell has a fixed enlargement of conditional mass greater than $F$;
these enlargements have bounded overlap.

For this fixed $Q$, let $B_x$ be its troublesome source set when the
pin survives the near-pin deletion, and let $B_x=\varnothing$ otherwise.
Put $\rho_Q=\rho|_Q/\rho(Q)$ and
\[
 I_Q=\int\rho_Q(B_x)\,d\lambda(x).
\]
We claim that $I_Q\le R^{-c_2}$, uniformly in $Q$. Suppose instead that
$I_Q\ge R^{-c_2}$. Jensen's inequality, applied to the pin-incidence
mass at each source, gives
\begin{equation}\label{ang:common-k-incidence}
 \int \rho_Q\Bigl(\bigcap_{i=1}^k B_{x_i}\Bigr)
           \,d\lambda(x_1)\cdots d\lambda(x_k)
      \ge I_Q^k.
\end{equation}
Tuples containing a removed pin contribute zero to this integral. For a
tuple of surviving pins, if
$|\det(v(x_1,Q),\ldots,v(x_k,Q))|<R^{-\eta}$, one of the
successive Gram--Schmidt distances is less than
$R^{-\eta/(k-1)}$. Conditional on the preceding pins, the corresponding
pin lies in a $C R^{-\eta/(k-1)}$-neighborhood of an affine hyperplane
through $z_Q$. Thus the product measure of these tuples is at most
\[
 C_kR^{\epsilon-\eta\zeta/(k-1)}.
\]
Choose $\eta$ small compared with $e'$, and then choose
$c_2,\epsilon$ so small that this is $o(R^{-kc_2})$. There is a tuple
of determinant at least $R^{-\eta}$ whose common source set has
conditional mass at least $\tfrac12R^{-kc_2}$. By
Equation~\eqref{ang:conditional-child}, it needs at least
\[
 c D R^{-kc_2-C_k\tau}
\]
balls of radius comparable to $2^{-b}$. On the other hand,
Lemma~\ref{ang:loomis-whitney}, applied to its $k$ tube covers, bounds
this number by
\[
 C D R^{k\eta-e'k/(k-1)}.
\]
The two estimates are contradictory if
\[
 k\eta+kc_2+C_k\tau<\frac{ke'}{k-1}.
\]
These parameter requirements are compatible: first choose $\eta$,
then $c_2$, and finally $\tau,\epsilon$. This proves the claim.
Deleting troublesome pairs and summing over $Q$ costs at most
$R^{-c_2}$ per block, because the weights $\rho(Q)$ sum to one.

Apply the endpoint density pruning of
Lemma~\ref{reg:pruning-multiplication}, with an additional loss
$R^{e''}$ per block, to the surviving pair set. Its physical-tube
version applies even though pins need not be separated from sources
by a fixed constant. To check this point directly, fix a parent $Q$ and
consider its depth-$b$ children meeting one final cap of radius
$2^{-N}$. Every point of such a child is within $C2^{-b}$ of the
target line. Choose one retained witness in one hit child of $Q$.
Its line differs from the target line by transverse distance
$O(2^{-N})$ in the bounded source region. Thus the pre-density-deletion
survivors in all hit children of this parent satisfy the same local
test at width $C'2^{-b}$. In each child their mass is at least
$R^{-e''}$ times the full child mass. The one local fraction bound
therefore charges the sum of the full masses of hit children by
$R^{e''}F$ times the full parent mass. It uses the
survivor set before the final density deletion; no renewed density
claim after that deletion is needed.

Multiplying these inequalities telescopes the profile increments.
Sum the resulting bounds over the occupied initial cells at the end of
the two ignored blocks. Their masses sum to one, so this introduces no
factor equal to the number of cells. The ignored growth is at most
$2k/J+o(1)$. Consequently every final cap has normalized class mass at most
\[
 R^{J(e'+e'')+o(1)}
 R^{-\frac{k-1}{k}(k-1-2k/J-C_k\tau-\epsilon)}.
\]
Choose $J$ large depending on $(k-1)^2/k-\ell_1$, then $e',e''$
small, and then the remaining losses as above. This proves the
single-scale assertion. All deletion exponents are positive fixed
numbers. Recombining classes with their original weights preserves
the cap estimate and the power deletion; the negligible discarded
class mass is included in the deletion.

Now choose $\ell<\ell_1<(k-1)^2/k$. It is enough to consider
$0<e<1$. Apply the single-scale assertion at dyadic $u$ between
$\delta$ and $\delta^{e/(2k)}$. At these scales the input loss
$\delta^{-\epsilon}$ is at most $u^{-2k\epsilon/e}$, so choosing
$\epsilon$ sufficiently small makes all the single-scale hypotheses
uniformly valid. The sum of the deletions is at most
\[
 C\log(1/\delta)\,\delta^{ec_1/(2k)},
\]
which is a positive power of $\delta$ after decreasing its exponent.
Intersect the surviving pair sets. Dyadic comparison gives the desired
cap estimate for $r\le\delta^{e/(2k)}$; the slack between $\ell$
and $\ell_1$ absorbs fixed constants. For larger $r$,
$\delta^{-e}r^\ell\ge1$ at small scales, so the estimate is trivial.
This completes the proof.
\end{proof}

\subsection{Proof of the linear gain}

\begin{proof}[Proof of Theorem~\ref{ang:linear-gain}]
We rule out any sequence with subpower input losses, subpower incidence
mass from below, and cover size $\delta^{-q-o(1)}$. This implies the
stated fixed-power result: if no positive $\epsilon,c$ worked, choose
counterexamples with $\epsilon=c=1/j$ at scales
$\delta_j<1/j$ tending to zero. They would form precisely such a
sequence. This observation will also justify all uses of the gain with
sufficiently small fixed losses.

By choosing one of finitely many small projective charts, we can assume
all relevant directions lie in a single chart and that incidence still
has subpower mass. If $q>k-1$, apply
Lemma~\ref{ang:slice-selection} with $h=k-1$. Select a slice with
subpower incidence and normalize its marginal restrictions; their
masses are bounded above by subpower and below by the incidence divided
by the other marginal mass. The coordinate maps take values in a plane, with
\[
 b=q-h\in(0,\tfrac12],\qquad
 f=t-h\in(\tfrac12,1),\qquad b<f,
\]
and tube count at most $\delta^{-b-o(1)}$. Apply
Lemma~\ref{ang:planar-measures} with the original sliced points and
directions as labels, and with coordinate maps $p_{V,z}$ and $v_V$.
Item~3 of Lemma~\ref{ang:slice-selection} gives the ball bounds on their
coordinate pushforwards and the separate cover for each original direction
label. The labeled incidence relation still has subpower mass, so these
data contradict the planar obstruction.

At $q=k-1$, slicing to the plane would leave direction exponent zero,
outside the range of Lemma~\ref{ang:planar-measures}. We therefore put
$h=k-2$, so the selected slices have dimension three, and let
\[
 g=t-h-\tfrac32>0.
\]
Choose once and for all $0<a<g/4$. We will find in one slice three
directions of determinant at least $c\delta^a$ with a common incident
point set of mass at least subpower. The resulting loss in
Lemma~\ref{ang:loomis-whitney} is $O(\delta^{-3a})$, which the
strict gap $g$ will absorb.

Parameterize our chart by $p\in\R^k$, with unit direction
$w(p)=(1,p)/\sqrt{1+|p|^2}$. Restricting and normalizing the chart
costs only a subpower factor; denote the resulting law in $p$ by
$\lambda$. In the rest of the endpoint proof the direction law supplied
to slicing is $w_\#\lambda$, and the incidence relation is restricted to
this chart. The slice law of the chart labels is
\[
 d\lambda_V(p)=\delta^{-h}
 \one_{\{|P_{V^\perp}w(p)|\le\delta\}}\,d\lambda(p),
 \qquad \nu_V=w_\#\lambda_V.
\]
Equation~\eqref{ang:endpoint-plate}
becomes the affine-hyperplane bound in the $p$ coordinates, because
an affine hyperplane there is the intersection of a linear hyperplane
with the chart. The ball exponent remains $k-1$. Choose
\[
 h<\ell<\frac{(k-1)^2}{k}=h+\frac1k.
\]
Apply Lemma~\ref{ang:weak-radial} to this chart law, with a small
fixed output loss $\delta^{-e}$. Its hypotheses hold at sufficiently
small scales for every fixed input-loss threshold. This deletes a
pair set of original product mass $O(\delta^{c_e})$ and gives, for
every anchor $p$, an $\ell$-cap bound $\delta^{-e}r^\ell$ on
the undeleted displacement directions. Also delete pairs at distance
less than $\delta^{b_0}$, where $b_0>0$ will be small. Denote the union
of these two exceptional pair relations by $\mathcal E_{\mathrm{pair}}$,
and call an anchored pair retained when it lies outside this union.

We first record how these deletions transfer to normalized slice-pair
mass. Equation~\eqref{ang:direction-pair-kernel}
shows that for every measurable chart-pair relation $\mathcal E$ with
$(\lambda\otimes\lambda)(\mathcal E)\le\delta^c$, and every fixed
$\theta>0$,
\begin{equation}\label{ang:exceptional-pair-transfer}
 \mathbb E_V(\lambda_V\otimes\lambda_V)(\mathcal E)
 \le C\delta^{c-\theta h}
       +\delta^{\theta(q-h)-o(1)}.
\end{equation}
For $0<\theta\le1$, the first term treats distances greater than
$\delta^\theta$, and the second treats smaller distances by the ball
bound and annular summation. For $\theta\ge1$, the kernel is at most
$C\delta^{-h}$, so the whole expectation is at most
$C\delta^{c-h}\le C\delta^{c-\theta h}$. The same annulus calculation
bounds the expected mass of the close pairs by
$\delta^{b_0(q-h)-o(1)}$. We leave $\theta$ free until $e$, and hence
the weak-radial deletion exponent $c_e$, has been fixed.

We next bound triples that are close to linearly dependent. For two
retained anchored pairs $(p,p'),(p,p'')$, set
\[
 \alpha=\left|
 \frac{p'-p}{|p'-p|}\wedge\frac{p''-p}{|p''-p|}\right|.
\]
In a bounded chart,
\begin{equation}\label{ang:chart-wedge}
 |w(p)\wedge w(p')\wedge w(p'')|
       \asymp |p'-p|\,|p''-p|\,\alpha.
\end{equation}
To see this, put $u=p'-p$ and $v=p''-p$ and subtract the first
unnormalized column from the other two. The orthogonal decomposition
of exterior powers gives the identity
\[
 |(1,p)\wedge(0,u)\wedge(0,v)|^2
       =|u\wedge v|^2+|p\wedge u\wedge v|^2.
\]
It supplies the lower bound, while boundedness of $p$ supplies the
upper bound. Normalizing the three columns changes only fixed constants.
Consequently a wedge smaller than $2\delta^a$ implies
$\alpha\le C\delta^{a-2b_0}$.

For $(w_1,w_2,w_3)=(w(p),w(p'),w(p''))$, the largest singular
value is comparable to one. Their first two columns span area at least $c\delta^{b_0}$,
so the second singular value is at least $c\delta^{b_0}$; by
Equation~\eqref{ang:chart-wedge}, the third is at least
$c\delta^{2b_0}\alpha$. Lemma~\ref{ang:haar} therefore bounds
their normalized expected slicing weight by
\begin{equation}\label{ang:triple-weight}
 \delta^{-3h}\Pr_V\{\max_{1\le i\le3}|P_{V^\perp}w_i|\le\delta\}
 \le C\delta^{-3hb_0}(\alpha+\delta)^{-h}.
\end{equation}
This includes singular values below $\delta$, since the kernel is
truncated there. At fixed $p,p'$, the weak radial estimate and
the two antipodal caps give
\[
 \lambda\{p'':(p,p'')\text{ is retained},\ \alpha\le r\}
       \le C\delta^{-e}r^\ell
       \qquad(\delta\le r\le1).
\]
Let $\mathcal T_\delta$ be the chart triples whose two anchored pairs
are retained and whose original wedge has size less than $2\delta^a$.
Summing dyadic annuli in Equation~\eqref{ang:triple-weight}, up to
$C\delta^{a-2b_0}$, gives
\begin{equation}\label{ang:bad-triples}
 \mathbb E_V\lambda_V^{\otimes3}(\mathcal T_\delta)
 \le C\delta^{(a-2b_0)(\ell-h)-e-3hb_0}.
\end{equation}
Since $\ell>h$, first choose $b_0>0$ sufficiently small depending
on $a,\ell,h$, and then $e>0$ sufficiently small, so that this
exponent is positive. These choices also ensure $a-2b_0>0$.
With $e$ fixed, the weak radial estimate fixes $c_e>0$. Now choose
$\theta>0$ small enough that $c_e-\theta h>0$. Since $q-h=1$, the
second exponent in Equation~\eqref{ang:exceptional-pair-transfer} is
also positive. Together with the close-pair bound, this proves that
$\mathcal E_{\mathrm{pair}}$ has power-small expected slice-pair mass.

Perform the restrictions in Lemma~\ref{ang:slice-selection}, and
write $\widetilde\lambda_V=(w^{-1})_\#\widetilde\nu_V$ for the restricted
law of chart labels. It is dominated by $\lambda_V$ and has the same
uniform subpower mass bound as $\widetilde\nu_V$. Define
\[
 F(V,z,x)=\widetilde\lambda_V
 \{p:(x,w(p))\in\widetilde E_{V,z}\}.
\]
Its integrated first moment is the retained sliced incidence and is at
least subpower, while the total integrated point mass is bounded by a
fixed constant.
H\"older's inequality consequently yields
\begin{equation}\label{ang:triple-incidence}
 \mathbb E_V\int\int F(V,z,x)^3
                 \,d\widetilde\mu_{V,z}(x)\,dz
       \ge\delta^{o(1)}.
\end{equation}
Dropping the point-incidence indicators bounds the contributions of
bad triples and exceptional anchored pairs by the left-hand sides below.
Equation~\eqref{ang:point-slice-total} and
$\widetilde\lambda_V\le\lambda_V$ give
\begin{align*}
 \mathbb E_V\int \|\widetilde\mu_{V,z}\|
      \widetilde\lambda_V^{\otimes3}(\mathcal T_\delta)\,dz
 &\le C\mathbb E_V\lambda_V^{\otimes3}(\mathcal T_\delta),\\
 \mathbb E_V\int \|\widetilde\mu_{V,z}\|\,\|\widetilde\lambda_V\|
      (\widetilde\lambda_V\otimes\widetilde\lambda_V)
          (\mathcal E_{\mathrm{pair}})\,dz
 &\le C\sup_V\|\widetilde\lambda_V\|\,
      \mathbb E_V(\lambda_V\otimes\lambda_V)
          (\mathcal E_{\mathrm{pair}}).
\end{align*}
Both quantities are power-small by Equation~\eqref{ang:bad-triples}
and the preceding bound for $\mathcal E_{\mathrm{pair}}$; the latter
also covers the two possible anchored pairs after changing the constant. No independence
between the point and direction slices is used. Hence a
subpower amount of the triple incidence in
Equation~\eqref{ang:triple-incidence} has original wedge at least
$2\delta^a$. Projection to $V$ and unit normalization change each
vector by $O(\delta)$, so their determinant in $V$ is at least
$\delta^a$ for small $\delta$.

The total parameter mass for integrating $(V,z)$ and three restricted
chart labels is at most subpower, since $\|\widetilde\lambda_V\|$ has a
uniform subpower bound and the $z$-domain is bounded.
It follows that some slice and some such triple have a common
incident point set of sliced mass at least subpower. Let $A$ be its
projected point coordinates. The pushforward of its restricted point
measure has the same mass and inherits the point-coordinate ball bound,
which implies
\[
 |A|_\delta\ge\delta^{-(t-h)+o(1)}.
\]
For each of the three original direction labels, the projected common
point set has its own cover by
$\delta^{-(q-h)-o(1)}=\delta^{-1-o(1)}$ tubes in that label's projected
direction. By
Lemma~\ref{ang:loomis-whitney},
\[
 |A|_\delta\le\delta^{-3/2-3a-o(1)}.
\]
This contradicts $3a<g=t-h-3/2$. When $k=2$ we have $h=0$;
all slicing weights equal one and the same proof applies, with the
$z$-integral interpreted as integration over the singleton
zero-dimensional space. The endpoint and the theorem are proved.
\end{proof}

\begin{corollary}[Small mass of bad directions]\label{ang:bad-directions}
Fix the parameters of Theorem~\ref{ang:linear-gain}. There are
$c_1,\epsilon_1>0$ such that the following holds at sufficiently small
$\delta$. Let $\mu$ be a bounded point probability satisfying the
$t$-ball bound with loss $\delta^{-\epsilon_1}$. Let $B_\delta$
be the set of directions for which at most
$\delta^{-q-c_1}$ width-$C\delta$ tubes capture $\mu$-mass at
least $\delta^{c_1}$. There is a Borel set
$B_\delta^+\supset B_\delta$, depending only on the point law and this
geometric test, such that every direction subprobability $\nu$
satisfying the $q$-ball bound with loss $\delta^{-\epsilon_1}$,
and, when $q=k-1$, Equation~\eqref{ang:endpoint-plate} with
$\epsilon_1$ in place of $\epsilon$, satisfies
\[
 \nu(B_\delta^+)\le\delta^{c_1}.
\]
In particular this bounds the outer $\nu$-measure of the original
geometric bad set. The direction law need not be normalized on either
set or on any individual incidence fiber. All predetermined fixed
widths are allowed.
\end{corollary}

\begin{proof}
Choose a fixed enlarged width $C'>C$ large enough for the spatial
anchor rounding used in Lemma~\ref{ang:slice-selection}. Let
$\epsilon,c$ be supplied by Theorem~\ref{ang:linear-gain} for width
$C'\delta$, and choose $\epsilon_1\le\epsilon/2$ and
$0<c_1<\min(\epsilon/2,c/3)$. In a fixed bounded spatial
$\delta$-grid consider all candidate lists $J$ of at most
$\lfloor\delta^{-q-c_1}\rfloor$ anchors. Keep the actual direction
$w$ in every tube, and define
\[
 B_\delta^+
 =\left\{w:\max_J
       \mu\left(\bigcup_{a\in J}T_{C'}(w,a)\right)
                    \ge\delta^{c_1}\right\}.
\]
The candidate masses are measurable, so this is a Borel set.
Rounding a point on each axis in an original bad cover gives a
candidate at the enlarged width with at least the same captured mass.
Consequently $B_\delta\subset B_\delta^+$; equality is not needed.
For each $w\in B_\delta^+$ choose the first candidate attaining the
threshold. Its measurable incidence relation has mass at least
$\delta^{c_1}\nu(B_\delta^+)$ and uses at most
$\delta^{-q-c_1}\le\delta^{-q-\epsilon}$ tubes per fiber.
Theorem~\ref{ang:linear-gain} bounds its incidence by $\delta^c$, hence
$\nu(B_\delta^+)\le\delta^{c-c_1}\le\delta^{c_1}$.
\end{proof}
